\newpage
\phantomsection
\label{prf:positive-fraction-iff-positive-product}

\begin{remark*}[Return]
\hyperref[thm:positive-fraction-iff-positive-product]{Return to Theorem}
\end{remark*}

\begin{theorem*}[Positivity of a Quotient via the Product $ab$]
For any $a,b\in F$ with $b\neq 0$,
\[
0<\frac{a}{b}
\quad\Longleftrightarrow\quad
0<a\cdot b.
\]
\end{theorem*}

\begin{proof}
\textbf{Professional Standard Proof.}~\\
Let $a,b\in F$ with $b\neq 0$.

\emph{($\Rightarrow$)} Assume
\[
0<\frac{a}{b}.
\]
By trichotomy, since $b\neq 0$, exactly one of the two possibilities
\[
b>0
\qquad\text{or}\qquad
b<0
\]
holds.

If $b>0$, then by
\hyperref[thm:multiply-positive]{Theorem~\ref*{thm:multiply-positive}},
multiplying the inequality
\[
0<\frac{a}{b}
\]
by the positive element $b$ gives
\[
0\cdot b<\frac{a}{b}\cdot b.
\]
Hence
\[
0<a.
\]
Multiplying again by the positive element $b$, we obtain
\[
0<a\cdot b.
\]

If $b<0$, then by
\hyperref[thm:multiply-negative]{Theorem~\ref*{thm:multiply-negative}},
multiplying the inequality
\[
0<\frac{a}{b}
\]
by the negative element $b$ reverses the order:
\[
0\cdot b>\frac{a}{b}\cdot b.
\]
Thus
\[
0>a,
\]
so $a<0$.  Since also $b<0$, the product of two negative elements is
positive.  Therefore
\[
0<a\cdot b.
\]

In either case,
\[
0<a\cdot b.
\]

\emph{($\Leftarrow$)} Assume
\[
0<a\cdot b.
\]
Again, since $b\neq 0$, either $b>0$ or $b<0$.

If $b>0$, then
\hyperref[thm:inverse-positive]{Theorem~\ref*{thm:inverse-positive}}
gives
\[
b^{-1}>0.
\]
Applying
\hyperref[thm:multiply-positive]{Theorem~\ref*{thm:multiply-positive}}
to
\[
0<a\cdot b
\]
with positive multiplier $b^{-1}$ yields
\[
0<(a\cdot b)b^{-1}.
\]
By associativity,
\[
(a\cdot b)b^{-1}=a(bb^{-1})=a.
\]
Hence
\[
0<a.
\]
Multiplying by the positive element $b^{-1}$ once more, we obtain
\[
0<ab^{-1}=\frac{a}{b}.
\]

If $b<0$, then $-b>0$, so by
\hyperref[thm:inverse-positive]{Theorem~\ref*{thm:inverse-positive}},
\[
(-b)^{-1}>0.
\]
Applying
\hyperref[thm:multiply-positive]{Theorem~\ref*{thm:multiply-positive}}
to
\[
0<a\cdot b
\]
with positive multiplier $(-b)^{-1}$ gives
\[
0<(a\cdot b)(-b)^{-1}.
\]
Since
\[
b(-b)^{-1}=-1,
\]
we have
\[
(a\cdot b)(-b)^{-1}=a\bigl(b(-b)^{-1}\bigr)=a(-1)=-a.
\]
Thus
\[
0<-a,
\]
so $a<0$.

Now both $a$ and $b$ are negative.  We compare $\frac{a}{b}$ and
$\frac{-a}{-b}$ using the definition of division.  Since
\[
b\cdot\frac{a}{b}=a,
\]
the negation-of-product identity gives
\[
(-b)\cdot\frac{a}{b}
=
-\left(b\cdot\frac{a}{b}\right)
=
-a.
\]
Thus $\frac{a}{b}$ is an element whose product with $-b$ is $-a$.  By the
definition of division, the unique such element is $\frac{-a}{-b}$.  Hence
\[
\frac{a}{b}=\frac{-a}{-b}.
\]
Since $-a>0$ and $-b>0$, the first case already proved implies
\[
0<\frac{-a}{-b}.
\]
Therefore
\[
0<\frac{a}{b}.
\]

Thus
\[
0<\frac{a}{b}
\quad\Longleftrightarrow\quad
0<a\cdot b.
\]
\end{proof}

\begin{remark*}[Proof structure]
The proof splits according to the sign of the denominator.  For the forward
direction, a positive quotient is multiplied by the denominator.  If the
denominator is positive, positivity is preserved and one gets $a>0$; if the
denominator is negative, order reverses and one gets $a<0$.  In either case,
$a$ and $b$ have the same sign, so their product is positive.

For the reverse direction, start from $ab>0$ and again inspect the sign of
$b$.  If $b>0$, multiplying by the positive inverse $b^{-1}$ shows first that
$a>0$ and then that $a/b>0$.  If $b<0$, multiplying by the positive inverse
$(-b)^{-1}$ shows that $-a>0$, so both $-a$ and $-b$ are positive; the
positive-denominator case then applies to $(-a)/(-b)$, and the quotient is
identified with $a/b$ by a sign-change identity.
\end{remark*}

\begin{remark*}[Dependencies]~\\
\begin{itemize}
  \item \hyperref[def:ordered-field]{Definition~\ref*{def:ordered-field}}
  \item \hyperref[thm:inverse-positive]{Theorem~\ref*{thm:inverse-positive}}
  \item \hyperref[thm:multiply-positive]{Theorem~\ref*{thm:multiply-positive}}
  \item \hyperref[thm:multiply-negative]{Theorem~\ref*{thm:multiply-negative}}
  \item \hyperref[thm:negation-of-product]{Theorem~\ref*{thm:negation-of-product}}
  \item \hyperref[cor:product-of-negatives]{Corollary~\ref*{cor:product-of-negatives}}
\end{itemize}
\end{remark*}
